% ====================================================================== Task 1
\section{Task 1: Stride Pattern Prefetcher at L1D}

% ---------------------------------------------------------------- data for the plots
% Every stage on every trace (Table 1). Columns: base ips v0 v1 p1 p2 p3 ac fin
\pgfplotstableread{
g      base   ips    v0     v1     p1     p2     p3     ac     fin
Trace1 0.2322 0.2334 0.2316 0.2321 0.2325 0.2342 0.2342 0.2344 0.2345
Trace2 0.2511 0.2512 0.2497 0.2499 0.2525 0.2535 0.2531 0.2537 0.2545
Trace3 0.3821 0.3830 0.3843 0.3836 0.3877 0.3858 0.3856 0.3869 0.3870
}\tStageIPC
\pgfplotstableread{
g      base   ips    v0     v1     p1     p2     p3     ac     fin
Trace1 1.0000 1.0051 0.9975 0.9996 1.0012 1.0084 1.0086 1.0094 1.0098
Trace2 1.0000 1.0002 0.9943 0.9951 1.0055 1.0093 1.0079 1.0102 1.0133
Trace3 1.0000 1.0024 1.0059 1.0042 1.0147 1.0097 1.0094 1.0128 1.0128
}\tStageSpd
\pgfplotstableread{
g      base   ips    v0     v1     p1     p2     p3     ac     fin
Trace1 10.68  10.12  9.70   9.82   9.50   9.83   9.89   9.70   9.68
Trace2 13.53  13.44  13.45  13.48  13.11  13.16  13.19  13.07  12.94
Trace3 3.40   3.32   3.07   3.21   2.91   3.10   3.12   3.02   3.04
}\tStageMpki
\pgfplotstableread{
g      base   ips    v0     v1     p1     p2     p3     ac     fin
Trace1 nan    52.8   32.4   35.4   38.2   56.0   61.4   56.5   57.8
Trace2 nan    38.1   18.7   18.7   26.6   54.7   63.0   59.3   63.3
Trace3 nan    16.1   18.6   20.5   31.3   37.3   40.5   41.9   41.0
}\tStageAcc
\pgfplotstableread{
g      base   ips    v0     v1     p1     p2     p3     ac     fin
Trace1 0.00   5.24   9.15   8.03   11.06  7.96   7.37   9.17   9.37
Trace2 0.00   0.73   0.59   0.43   3.11   2.77   2.53   3.45   4.40
Trace3 0.00   2.51   9.83   5.64   14.48  8.96   8.25   11.13  10.69
}\tStageCov
\newcommand{\stagecols}{base,ips,v0,v1,p1,p2,p3,ac,fin}
\newcommand{\stageopts}{cycle list name=stages, bar width=2.5pt, enlarge x limits=0.2}

\subsection{Task 1A: IP-Stride Prefetcher (\code{IP-stride})}
ChampSim calls the L1D prefetcher on every demand load, with the virtual address and the instruction
pointer (IP) of the load. The provided \code{ip\_stride} prefetcher works like this:
\begin{itemize}
  \item \textbf{Tracker table:} 24 entries, fully associative, replaced by LRU. Each entry stores an IP tag
        (\code{ip \& 0xFF}), the last cache-line address of that IP, the last stride (in cache lines) and the LRU rank.
  \item \textbf{New IP:} the LRU entry is replaced and no prefetch is issued.
  \item \textbf{Known IP:} the new stride is (current line $-$ last line). A stride of 0 is ignored.
  \item \textbf{Trigger:} a prefetch is issued only when the new stride equals the stored stride (the same stride twice in a row).
  \item \textbf{Prefetch:} lines $\mathrm{cl}+s$, $\mathrm{cl}+2s$, $\mathrm{cl}+3s$ (degree 3). The loop stops when a line
        leaves the 4\,KB page of the demand access.
  \item \textbf{Target level:} the L1D if less than half of the MSHRs are used, otherwise the L2.
  \item \textbf{Feedback:} none. The \code{cache\_fill} hook is empty and the degree is fixed.
\end{itemize}
It captures only constant strides such as $A$, $A{+}2$, $A{+}4$. A repeating pattern such as
$+1, +2, +2, +2, +3$ never has two equal strides in a row, so it never triggers a prefetch.

\subsection{Task 1B: Stride Pattern Prefetcher (SPP)}
SPP runs at the L1D, tracks every IP separately, never crosses a 4\,KB page, supports a degree larger than 1,
and checks \emph{uniform} strides first and \emph{non-uniform} stride patterns second. The sections below add one
feature at a time. The names \code{SPP-V0}, \code{SPP-V1}, \code{SPP-P1}, \code{SPP-P2}, \code{SPP-P3}, \code{SPP-AC}
and \code{SPP-Final} are the names used in the tables and plots.

\subsubsection{\code{SPP-V0}: Basic Non-Uniform Prefetcher}
\begin{itemize}
  \item \textbf{Tracker table:} same as \code{IP-stride} (24 entries, 8-bit IP tag, LRU), with two extra fields per
        entry: a 12-bit \emph{signature} and a \emph{uniform counter}.
  \item \textbf{Signature:} the last three strides of the IP, 4 bits each:
        \[ \mathrm{sig}' = \bigl((\mathrm{sig} \ll 4) \mid (s \mathbin{\&} \code{0xF})\bigr) \mathbin{\&} \code{0xFFF} \]
  \item \textbf{Pattern History Table (PHT):} 4096 entries, indexed by the signature, each storing one stride.
        When a new stride $s$ is seen, the entry of the \emph{old} signature is set to $s$
        (``after this history, this stride came next'').
  \item \textbf{Uniform first:} if the new stride equals the previous stride, the IP is uniform and SPP prefetches
        exactly like \code{IP-stride}.
  \item \textbf{Non-uniform fallback:} otherwise the signature reads the PHT to get the next stride. The signature
        is then advanced with that predicted stride to predict the stride after it (chaining), until the degree
        is reached or the PHT entry is empty.
  \item \textbf{Degree:} fixed at 3. Prefetches stay in the demand page. Target level (L1D or L2) is chosen by MSHR
        occupancy, as in \code{IP-stride}.
\end{itemize}
\emph{Example:} for the repeating strides $(+1, +2, +2, +2, +3)$ the five 3-stride signatures are all different,
$(1,2,2)\to2$, $(2,2,2)\to3$, $(2,2,3)\to1$, $(2,3,1)\to2$, $(3,1,2)\to2$, so after one pass the PHT predicts the
whole pattern.

\subsubsection{\code{SPP-V1}: Fill-Based Throttling}
\code{SPP-V0} plus a first throttling step:
\begin{itemize}
  \item Every 1000 demand accesses, compute (prefetched lines filled into L1D) / (prefetches issued).
  \item Degree 3 if this ratio is $\ge 50\%$, degree 2 if $\ge 25\%$, otherwise degree 1.
  \item This ratio counts prefetches that arrived in the cache, not prefetches that were used later.
\end{itemize}

\subsubsection{\code{SPP-P1}: Larger Tracker Table}
\code{SPP-V1} plus:
\begin{itemize}
  \item Tracker table from 24 to 128 entries, IP tag from 8 to 12 bits.
  \item A valid bit per entry, so an empty entry cannot match an IP whose tag is 0.
  \item \textbf{Why:} fewer different IPs share one tracker, and trackers live long enough to collect a full stride history.
\end{itemize}

\subsubsection{\code{SPP-P2}: Clean Stride History}
\code{SPP-P1} plus:
\begin{itemize}
  \item \textbf{Stride range:} a stride outside $[-64, +63]$ lines restarts the history of that IP and is not stored.
        A page has only 64 lines, so such a stride can never give a same-page prefetch.
  \item \textbf{History count:} a 2-bit counter (0 to 3) per tracker. The PHT is trained and read only after 3 valid
        strides, so signatures padded with zeros do not pollute it.
\end{itemize}

\subsubsection{\code{SPP-P3}: PHT Confidence}
\code{SPP-P2} plus a 2-bit saturating confidence in every PHT entry (0 means empty):
\begin{itemize}
  \item \textbf{Entry empty:} store the stride, confidence $=1$.
  \item \textbf{Same stride seen again:} confidence $+1$ (maximum 3).
  \item \textbf{Different stride seen:} confidence $-1$. At 0 the entry takes the new stride with confidence 1.
  \item \textbf{Prediction:} the first predicted stride needs confidence $\ge 2$. Every further (chained) stride needs
        confidence $=3$, because a chained prediction depends on the one before it.
\end{itemize}

\subsubsection{\code{SPP-AC}: Accuracy Checker}
\code{SPP-AC} is \code{SPP-P3} with the fill-based throttling of \code{SPP-V1} replaced by a checker that measures real usefulness.
\begin{itemize}
  \item \textbf{Prefetch filter:} 64 entries, direct-mapped by the low 6 bits of the cache-line address, each with a
        valid bit and a 10-bit partial tag. It remembers lines our prefetcher filled into the L1D that no demand
        load has used yet.
  \item \textbf{Fill hook:} when a prefetched line is filled (\code{l1d\_prefetcher\_cache\_fill}, \code{prefetch=1}),
        insert its virtual line address into the filter and add 1 to the fill counter $F$.
  \item \textbf{Use hook:} a demand load that hits in the L1D (\code{l1d\_prefetcher\_operate}) and finds its line in
        the filter adds 1 to the useful counter $U$ and clears the entry, so each line counts once.
  \item \textbf{Epoch:} after $W = 256$ fills, the accuracy is $A = \lfloor 100\,U/F \rfloor$ and the degree for the next epoch is
        \[ d = \begin{cases}
            D_H & A \ge T_H\\
            D_M & T_M \le A < T_H\\
            D_L & T_L \le A < T_M\\
            D_{\min} & A < T_L
          \end{cases} \]
        Then $F$ and $U$ are reset to 0. The degree starts at 3.
  \item $\boldsymbol{D_{\min} = 1}$\textbf{, never 0:} epochs are counted in fills, so with degree 0 no fill would happen,
        the epoch would never end and the degree could never go up again.
\end{itemize}
\code{SPP-AC} uses 128 trackers with a 12-bit IP tag, a 64-entry filter, $W = 256$, thresholds
$(T_H, T_M, T_L) = (75, 50, 25)$ and degrees $(D_H, D_M, D_L, D_{\min}) = (6, 4, 2, 1)$. Section~\ref{sec:t1-values}
examines each of these values.

\subsubsection{\code{SPP-Final}: The Submitted Prefetcher}
\code{SPP-Final} is \code{SPP-AC} with an \textbf{8-bit IP tag} and a \textbf{32-entry prefetch filter}. All other values
are the same as in \code{SPP-AC}. It is the file \code{prefetcher/spp.l1d\_pref}. Section~\ref{sec:t1-final} explains why.

\subsubsection{Sweep Variants}
Each variant changes one value of \code{SPP-AC}.
\begin{table}[H]
\centering\small
\begin{tabular}{@{}ll@{}}
\toprule
Names & What changes \\
\midrule
\code{Trk64}, \code{Trk256}         & tracker entries 64 or 256 (default 128) \\
\code{Tag8}                         & 8-bit IP tag (default 12) \\
\code{Flt32}, \code{Flt128}         & prefetch filter 32 or 128 entries (default 64) \\
\code{Win128}, \code{Win512}        & epoch of 128 or 512 fills (default 256) \\
\code{ThrStrict}, \code{ThrLoose}   & thresholds 90/70/50 or 60/40/20 (default 75/50/25) \\
\code{DegHi4}, \code{DegHi8}        & high-band degree 4 or 8 (default 6) \\
\code{DegMid3}, \code{DegMid5}      & middle-band degree 3 or 5 (default 4) \\
\code{DegLow1}, \code{DegLow3}      & low-band degree 1 or 3 (default 2) \\
\code{Fix1} to \code{Fix4}          & throttling off, always degree 1 to 4 \\
\bottomrule
\end{tabular}
\end{table}

\subsection{Metrics}
$I$ is the number of simulated instructions, $C$ the number of cycles and $M$ the number of L1D load misses.
Subscript $b$ is the baseline (no prefetching), $p$ is a run with a prefetcher.
\[
\mathrm{IPC} = I/C,\qquad \mathrm{Speedup} = \mathrm{IPC}_p/\mathrm{IPC}_b,\qquad \mathrm{MPKI} = 1000\,M/I,
\]
\[
\mathrm{Coverage} = (M_b - M_p)/M_b \times 100\%,\qquad
\mathrm{Accuracy} = \frac{\text{useful prefetches}}{\text{prefetches issued to the next level}} \times 100\%.
\]
MPKI and accuracy are read from the ChampSim statistics of the L1D. The accuracy checker inside \code{SPP-AC} and
\code{SPP-Final} uses its own cheaper estimate only for throttling, never in the tables. When we compare
configurations over all traces we use the geometric mean of the three speedups, and the arithmetic mean of
accuracy and coverage. All runs: single core, hashed perceptron branch predictor, LRU at all levels, no other
prefetcher, 25M warm-up and 25M simulated instructions, traces Trace1 to Trace3.

\subsection{Impact of Each Feature}
Table~\ref{tab:t1-stages} shows every stage on every trace, Table~\ref{tab:t1-avg} the averages, and
Figures~\ref{fig:t1-stages} and~\ref{fig:t1-cov} the same data as plots. Each stage adds one feature to the stage
before it, in the order of Section~1.2.

\begin{table}[H]
\centering\small
\begin{tabular}{@{}lccccc@{}}
\toprule
Prefetcher & IPC & Speedup & L1D load MPKI & Accuracy (\%) & Coverage (\%) \\
\midrule
\grp{6}{Trace1}
Baseline           & 0.2322 & 1.0000 & 10.68 & \na  & 0.00 \\
\code{IP-stride}   & 0.2334 & 1.0051 & 10.12 & 52.8 & 5.24 \\
\code{SPP-V0}      & 0.2316 & 0.9975 & 9.70  & 32.4 & 9.15 \\
\code{SPP-V1}      & 0.2321 & 0.9996 & 9.82  & 35.4 & 8.03 \\
\code{SPP-P1}      & 0.2325 & 1.0012 & 9.50  & 38.2 & 11.06 \\
\code{SPP-P2}      & 0.2342 & 1.0084 & 9.83  & 56.0 & 7.96 \\
\code{SPP-P3}      & 0.2342 & 1.0086 & 9.89  & 61.4 & 7.37 \\
\code{SPP-AC}      & 0.2344 & 1.0094 & 9.70  & 56.5 & 9.17 \\
\code{SPP-Final}   & 0.2345 & 1.0098 & 9.68  & 57.8 & 9.37 \\
\midrule
\grp{6}{Trace2}
Baseline           & 0.2511 & 1.0000 & 13.53 & \na  & 0.00 \\
\code{IP-stride}   & 0.2512 & 1.0002 & 13.44 & 38.1 & 0.73 \\
\code{SPP-V0}      & 0.2497 & 0.9943 & 13.45 & 18.7 & 0.59 \\
\code{SPP-V1}      & 0.2499 & 0.9951 & 13.48 & 18.7 & 0.43 \\
\code{SPP-P1}      & 0.2525 & 1.0055 & 13.11 & 26.6 & 3.11 \\
\code{SPP-P2}      & 0.2535 & 1.0093 & 13.16 & 54.7 & 2.77 \\
\code{SPP-P3}      & 0.2531 & 1.0079 & 13.19 & 63.0 & 2.53 \\
\code{SPP-AC}      & 0.2537 & 1.0102 & 13.07 & 59.3 & 3.45 \\
\code{SPP-Final}   & 0.2545 & 1.0133 & 12.94 & 63.3 & 4.40 \\
\midrule
\grp{6}{Trace3}
Baseline           & 0.3821 & 1.0000 & 3.40 & \na  & 0.00 \\
\code{IP-stride}   & 0.3830 & 1.0024 & 3.32 & 16.1 & 2.51 \\
\code{SPP-V0}      & 0.3843 & 1.0059 & 3.07 & 18.6 & 9.83 \\
\code{SPP-V1}      & 0.3836 & 1.0042 & 3.21 & 20.5 & 5.64 \\
\code{SPP-P1}      & 0.3877 & 1.0147 & 2.91 & 31.3 & 14.48 \\
\code{SPP-P2}      & 0.3858 & 1.0097 & 3.10 & 37.3 & 8.96 \\
\code{SPP-P3}      & 0.3856 & 1.0094 & 3.12 & 40.5 & 8.25 \\
\code{SPP-AC}      & 0.3869 & 1.0128 & 3.02 & 41.9 & 11.13 \\
\code{SPP-Final}   & 0.3870 & 1.0128 & 3.04 & 41.0 & 10.69 \\
\bottomrule
\end{tabular}
\caption{Every stage of the design on every trace (25M warm-up, 25M simulated instructions).}
\label{tab:t1-stages}
\end{table}

\begin{table}[H]
\centering\small
\begin{tabular}{@{}lccc@{}}
\toprule
Prefetcher & Geometric-mean speedup & Mean accuracy (\%) & Mean coverage (\%) \\
\midrule
Baseline         & 1.00000 & \na   & 0.00 \\
\code{IP-stride} & 1.00257 & 35.67 & 2.83 \\
\code{SPP-V0}    & 0.99923 & 23.21 & 6.52 \\
\code{SPP-V1}    & 0.99963 & 24.88 & 4.70 \\
\code{SPP-P1}    & 1.00710 & 32.04 & 9.55 \\
\code{SPP-P2}    & 1.00914 & 49.30 & 6.56 \\
\code{SPP-P3}    & 1.00863 & 55.00 & 6.05 \\
\code{SPP-AC}    & 1.01079 & 52.58 & 7.92 \\
\code{SPP-Final} & 1.01200 & 54.03 & 8.15 \\
\bottomrule
\end{tabular}
\caption{Average over the three traces (geometric mean for speedup, arithmetic mean for accuracy and coverage).}
\label{tab:t1-avg}
\end{table}

\begin{figure}[H]
\centering
\begin{subfigure}{0.49\textwidth}\centering
  \gbarplot{\tStageIPC}{\stagecols}{\stageopts, ymin=0, ymax=0.45, ylabel={IPC},
    legend entries={Baseline,IP-stride,V0,V1,P1,P2,P3,AC,Final}, legend columns=5,
    legend to name=leg:stages}
  \caption{IPC}
\end{subfigure}\hfill
\begin{subfigure}{0.49\textwidth}\centering
  \gbarplot{\tStageSpd}{\stagecols}{\stageopts, ymin=0.99, ymax=1.02, ylabel={Speedup}, speedline}
  \caption{Speedup}
\end{subfigure}\\[4pt]
\begin{subfigure}{0.49\textwidth}\centering
  \gbarplot{\tStageMpki}{\stagecols}{\stageopts, ymin=0, ymax=15, ylabel={L1D load MPKI}}
  \caption{L1D load MPKI}
\end{subfigure}\hfill
\begin{subfigure}{0.49\textwidth}\centering
  \gbarplot{\tStageAcc}{\stagecols}{\stageopts, ymin=0, ymax=70, ylabel={Accuracy (\%)}}
  \caption{Prefetch accuracy}
\end{subfigure}\\[3pt]
\ref*{leg:stages}
\caption{Every stage of the design. Bars in each trace group, left to right: Baseline, \code{IP-stride},
\code{SPP-V0}, \code{SPP-V1}, \code{SPP-P1}, \code{SPP-P2}, \code{SPP-P3}, \code{SPP-AC}, \code{SPP-Final}.}
\label{fig:t1-stages}
\end{figure}

\begin{figure}[H]
\centering
\begin{subfigure}{0.49\textwidth}\centering
  \gbarplot{\tStageCov}{\stagecols}{\stageopts, ymin=0, ymax=16, ylabel={Coverage (\%)}}
  \caption{Prefetch coverage}
\end{subfigure}\\[3pt]
\ref*{leg:stages}
\caption{Coverage of every stage (same bar order as Figure~\ref{fig:t1-stages}). The baseline issues no prefetches,
so it has no accuracy or coverage bars.}
\label{fig:t1-cov}
\end{figure}

\subsection{How Each Value Was Chosen}\label{sec:t1-values}

\subsubsection{Method}
\begin{itemize}
  \item \textbf{Starting point:} \code{SPP-AC}. Each sweep variant changes \textbf{one} value and keeps all others,
        with the same traces and instruction counts.
  \item \textbf{Main measure:} geometric-mean speedup over the three traces. Accuracy and coverage are used to explain the result.
  \item A gap below 0.05\% is called a \textbf{tie}. We only have three traces and cannot tell such a small gap from
        trace-sample noise. The 0.05\% cutoff is our own choice.
  \item In a tie we keep the current value. The exception is table size: equal speed with less storage wins.
  \item Each value was changed alone. The two changes that make \code{SPP-Final} were then run together (Section~\ref{sec:t1-final}).
  \item Every subsection below gives the table and the plots of one value, then the decision with its reason.
        The plots use the same axis range, so they can be compared.
\end{itemize}

% sweep tables: \swrow{name}{value}{T1}{T2}{T3}{geomean}{acc}{cov}
\newcommand{\swhead}{\toprule Variant & Value & Speedup T1 & Speedup T2 & Speedup T3 & Geomean speedup & Mean acc (\%) & Mean cov (\%) \\ \midrule}
\newcommand{\swtab}[2]{%
  \begin{table}[H]\centering\small
  \begin{tabular}{@{}llcccccc@{}}\swhead #1 \bottomrule\end{tabular}
  \caption{#2 Only this value changes; bold row = \code{SPP-AC}.}
  \end{table}}
\newcommand{\acrow}{\best\textbf{\code{SPP-AC}}}
\newcommand{\acvals}{1.0094 & 1.0102 & 1.0128 & 1.01079 & 52.58 & 7.92}

\subsubsection{IP Tracker Table Entries}
\code{Trk64}, \code{SPP-AC} (128) and \code{Trk256}.
\swtab{
  \code{Trk64}  & 64  & 1.0083 & 1.0026 & 1.0082 & 1.00636 & 47.51 & 6.16 \\
  \acrow        & 128 & \acvals \\
  \code{Trk256} & 256 & 1.0107 & 1.0124 & 1.0162 & 1.01313 & 53.89 & 8.98 \\
}{IP tracker table entries.}
\pgfplotstableread{
g      a      b      c
Trace1 1.0083 1.0094 1.0107
Trace2 1.0026 1.0102 1.0124
Trace3 1.0082 1.0128 1.0162
}\tSw
\pgfplotstableread{
g      acc   cov
Trk64  47.51 6.16
SPP-AC 52.58 7.92
Trk256 53.89 8.98
}\tAc
\sweepfig{\tSw}{a,b,c}{Trk64,SPP-AC,Trk256}{\tAc}{IP tracker table entries: speedup per trace (left) and mean accuracy and coverage (right).}
\begin{itemize}
  \item 64 entries is clearly worse ($-0.439\%$).
  \item 256 entries is better ($+0.232\%$) and better on all three traces.
  \item \textbf{Cost:} 64 more entries (64 to 128) need $64 \times 84 = 5376$ bits, about 0.66\,KiB, and gave about 0.67\%
        speedup per KiB. 128 more entries (128 to 256) need $128 \times 84 = 10752$ bits, about 1.3\,KiB (22\% of the
        whole prefetcher), and give about 0.18\% speedup per KiB. The second doubling gives under a third of the
        speedup per KiB of the first.
  \item \textbf{Decision: keep 128 entries.} The assignment asks for a lightweight prefetcher with maximum performance
        for minimal storage. The gain of the second doubling is real but small for its cost, so we accept the lower
        speed to stay small. If storage were not a concern, 256 would be the better choice.
\end{itemize}

\subsubsection{IP Tag Bits}
\code{Tag8} and \code{SPP-AC} (12 bits), with 128 entries.
\swtab{
  \code{Tag8} & 8  & 1.0098 & 1.0135 & 1.0130 & 1.01206 & 53.86 & 8.28 \\
  \acrow      & 12 & \acvals \\
}{IP tag bits.}
\pgfplotstableread{
g      a      b
Trace1 1.0098 1.0094
Trace2 1.0135 1.0102
Trace3 1.0130 1.0128
}\tSw
\pgfplotstableread{
g      acc   cov
Tag8   53.86 8.28
SPP-AC 52.58 7.92
}\tAc
\sweepfig{\tSw}{a,b}{Tag8,SPP-AC}{\tAc}{IP tag bits: speedup per trace (left) and mean accuracy and coverage (right).}
\begin{itemize}
  \item The 8-bit tag is faster ($+0.126\%$, a real gain) and also needs 4 bits less per tracker entry.
  \item \textbf{Decision: 8 bits.} Better and smaller.
  \item We tested only 8 and 12 bits, so 8 is not shown to be the best width. We did not investigate why the shorter tag helps.
\end{itemize}

\subsubsection{Prefetch Filter Entries}
\code{Flt32}, \code{SPP-AC} (64) and \code{Flt128}.
\swtab{
  \code{Flt32}  & 32  & 1.0094 & 1.0101 & 1.0128 & 1.01075 & 52.78 & 7.85 \\
  \acrow        & 64  & \acvals \\
  \code{Flt128} & 128 & 1.0094 & 1.0102 & 1.0131 & 1.01092 & 52.55 & 8.02 \\
}{Prefetch filter entries.}
\pgfplotstableread{
g      a      b      c
Trace1 1.0094 1.0094 1.0094
Trace2 1.0101 1.0102 1.0102
Trace3 1.0128 1.0128 1.0131
}\tSw
\pgfplotstableread{
g      acc   cov
Flt32  52.78 7.85
SPP-AC 52.58 7.92
Flt128 52.55 8.02
}\tAc
\sweepfig{\tSw}{a,b,c}{Flt32,SPP-AC,Flt128}{\tAc}{Prefetch filter entries: speedup per trace (left) and mean accuracy and coverage (right).}
\begin{itemize}
  \item 32, 64 and 128 entries are ties ($-0.004\%$ and $+0.013\%$). Accuracy and coverage barely change.
  \item \textbf{Decision: 32 entries.} Equal speed with less storage wins.
  \item We did not test fewer than 32 entries.
\end{itemize}

\subsubsection{Epoch Length}
\code{Win128}, \code{SPP-AC} (256) and \code{Win512}.
\swtab{
  \code{Win128} & 128 & 1.0095 & 1.0103 & 1.0134 & 1.01106 & 52.81 & 8.04 \\
  \acrow        & 256 & \acvals \\
  \code{Win512} & 512 & 1.0092 & 1.0102 & 1.0124 & 1.01058 & 51.40 & 7.60 \\
}{Epoch length (prefetch fills).}
\pgfplotstableread{
g      a      b      c
Trace1 1.0095 1.0094 1.0092
Trace2 1.0103 1.0102 1.0102
Trace3 1.0134 1.0128 1.0124
}\tSw
\pgfplotstableread{
g      acc   cov
Win128 52.81 8.04
SPP-AC 52.58 7.92
Win512 51.40 7.60
}\tAc
\sweepfig{\tSw}{a,b,c}{Win128,SPP-AC,Win512}{\tAc}{Epoch length (prefetch fills): speedup per trace (left) and mean accuracy and coverage (right).}
\begin{itemize}
  \item 128, 256 and 512 fills are ties ($+0.027\%$ and $-0.020\%$).
  \item \textbf{Decision: keep 256.} With 256 samples the measured accuracy has a standard error of at most
        $\sqrt{0.25/256} \approx 3.1$ points, much smaller than the 25-point band width, so a band is rarely chosen
        by noise. A shorter epoch would react faster but change the degree more often, a longer one would react
        slowly to program phases. The sweep shows no reason to change it.
\end{itemize}

\subsubsection{Accuracy Thresholds}
\code{ThrStrict} (90/70/50), \code{SPP-AC} (75/50/25) and \code{ThrLoose} (60/40/20), with the degrees 6/4/2/1.
\swtab{
  \code{ThrStrict} & 90/70/50 & 1.0092 & 1.0101 & 1.0116 & 1.01026 & 56.24 & 7.13 \\
  \acrow           & 75/50/25 & \acvals \\
  \code{ThrLoose}  & 60/40/20 & 1.0090 & 1.0102 & 1.0132 & 1.01080 & 50.53 & 8.15 \\
}{Accuracy thresholds.}
\pgfplotstableread{
g      a      b      c
Trace1 1.0092 1.0094 1.0090
Trace2 1.0101 1.0102 1.0102
Trace3 1.0116 1.0128 1.0132
}\tSw
\pgfplotstableread{
g         acc   cov
ThrStrict 56.24 7.13
SPP-AC    52.58 7.92
ThrLoose  50.53 8.15
}\tAc
\sweepfig{\tSw}{a,b,c}{ThrStrict,SPP-AC,ThrLoose}{\tAc}{Accuracy thresholds: speedup per trace (left) and mean accuracy and coverage (right).}
\begin{itemize}
  \item Starting point: four degrees need four accuracy bands, so we split 0 to 100\% into four equal bands, 75/50/25.
  \item \code{ThrStrict} is $-0.053\%$ against 75/50/25. Its accuracy is higher (56.2\% against 52.6\%) but its coverage
        is lower (7.13\% against 7.92\%): it keeps the degree low too often.
  \item \code{ThrLoose} is a tie in speedup ($+0.001\%$) and has lower accuracy (50.5\%).
  \item \textbf{Decision: keep 75/50/25.} Strict is slower, and loose is not faster and is less accurate.
\end{itemize}

\subsubsection{Prefetch Degrees: Why Different Degrees for Different Accuracy}
The degree is how far ahead we prefetch. Going further hides more latency when the prefetches are used, but every
extra prefetch that is not used wastes bandwidth and takes a place in the L1D. So accurate streams get a high
degree and inaccurate streams a low degree. The next four tables test this: fixed degrees first, then each band alone.

\begin{table}[H]\centering\small
\begin{tabular}{@{}llcccccc@{}}\swhead
  \code{Fix1} & 1 & 1.0081 & 1.0078 & 1.0075 & 1.00780 & 55.12 & 5.28 \\
  \code{Fix2} & 2 & 1.0091 & 1.0086 & 1.0102 & 1.00930 & 51.86 & 6.65 \\
  \code{Fix3} & 3 & 1.0091 & 1.0088 & 1.0115 & 1.00981 & 47.75 & 7.39 \\
  \code{Fix4} & 4 & 1.0091 & 1.0091 & 1.0125 & 1.01024 & 44.77 & 7.99 \\
\bottomrule\end{tabular}
\caption{Fixed degree, throttling off. Only this value changes; bold row = \code{SPP-AC}.}
\end{table}
\pgfplotstableread{
g      a      b      c      d
Trace1 1.0081 1.0091 1.0091 1.0091
Trace2 1.0078 1.0086 1.0088 1.0091
Trace3 1.0075 1.0102 1.0115 1.0125
}\tSw
\pgfplotstableread{
g    acc   cov
Fix1 55.12 5.28
Fix2 51.86 6.65
Fix3 47.75 7.39
Fix4 44.77 7.99
}\tAc
\sweepfig{\tSw}{a,b,c,d}{Fix1,Fix2,Fix3,Fix4}{\tAc}{Fixed degree, throttling off: speedup per trace (left) and mean accuracy and coverage (right).}

\textbf{Fixed degrees (throttling off)}
\begin{itemize}
  \item Speedup grows with the degree, but accuracy falls (55.1\%, 51.9\%, 47.8\%, 44.8\%). A deeper degree gives
        speed but costs accuracy.
  \item \code{SPP-AC} (throttled, 6/4/2/1) is $+0.054\%$ against the best fixed degree (\code{Fix4}), with accuracy
        52.6\% against 44.8\%. The checker gives high degrees only when prefetches are being used.
\end{itemize}

\textbf{High band (accuracy at least 75\%)}\\
\code{DegHi4}, \code{SPP-AC} (6) and \code{DegHi8}.
\swtab{
  \code{DegHi4} & 4 & 1.0090 & 1.0094 & 1.0116 & 1.01004 & 51.56 & 7.55 \\
  \acrow        & 6 & \acvals \\
  \code{DegHi8} & 8 & 1.0092 & 1.0109 & 1.0136 & 1.01125 & 53.33 & 8.18 \\
}{Degree of the high band.}
\pgfplotstableread{
g      a      b      c
Trace1 1.0090 1.0094 1.0092
Trace2 1.0094 1.0102 1.0109
Trace3 1.0116 1.0128 1.0136
}\tSw
\pgfplotstableread{
g      acc   cov
DegHi4 51.56 7.55
SPP-AC 52.58 7.92
DegHi8 53.33 8.18
}\tAc
\sweepfig{\tSw}{a,b,c}{DegHi4,SPP-AC,DegHi8}{\tAc}{Degree of the high band: speedup per trace (left) and mean accuracy and coverage (right).}
\begin{itemize}
  \item Degree 4 is worse ($-0.074\%$). Degree 8 is a tie ($+0.045\%$), better on Trace2 and Trace3 but lower on Trace1.
  \item \textbf{Decision: keep 6.} When at least 75\% of prefetches are used, a deeper degree pays off, as the
        fixed-degree test shows. 8 is not clearly better than 6, so we keep the smaller degree.
\end{itemize}

\textbf{Middle band (accuracy 50 to 75\%)}\\
\code{DegMid3}, \code{SPP-AC} (4) and \code{DegMid5}.
\swtab{
  \code{DegMid3} & 3 & 1.0094 & 1.0102 & 1.0125 & 1.01067 & 53.30 & 7.78 \\
  \acrow         & 4 & \acvals \\
  \code{DegMid5} & 5 & 1.0090 & 1.0103 & 1.0133 & 1.01085 & 52.29 & 8.04 \\
}{Degree of the middle band.}
\pgfplotstableread{
g      a      b      c
Trace1 1.0094 1.0094 1.0090
Trace2 1.0102 1.0102 1.0103
Trace3 1.0125 1.0128 1.0133
}\tSw
\pgfplotstableread{
g       acc   cov
DegMid3 53.30 7.78
SPP-AC  52.58 7.92
DegMid5 52.29 8.04
}\tAc
\sweepfig{\tSw}{a,b,c}{DegMid3,SPP-AC,DegMid5}{\tAc}{Degree of the middle band: speedup per trace (left) and mean accuracy and coverage (right).}
\begin{itemize}
  \item Degree 3 ($-0.011\%$) and degree 5 ($+0.006\%$) are ties.
  \item \textbf{Decision: keep 4}, between the high-band degree 6 and the low-band degree 2. The data does not show
        that 4 is better than 3 or 5. It is kept as chosen.
\end{itemize}

\textbf{Low band (accuracy 25 to 50\%)}\\
\code{DegLow1}, \code{SPP-AC} (2) and \code{DegLow3}.
\swtab{
  \code{DegLow1} & 1 & 1.0090 & 1.0100 & 1.0122 & 1.01042 & 54.31 & 7.43 \\
  \acrow         & 2 & \acvals \\
  \code{DegLow3} & 3 & 1.0093 & 1.0101 & 1.0132 & 1.01086 & 50.29 & 8.03 \\
}{Degree of the low band.}
\pgfplotstableread{
g      a      b      c
Trace1 1.0090 1.0094 1.0093
Trace2 1.0100 1.0102 1.0101
Trace3 1.0122 1.0128 1.0132
}\tSw
\pgfplotstableread{
g       acc   cov
DegLow1 54.31 7.43
SPP-AC  52.58 7.92
DegLow3 50.29 8.03
}\tAc
\sweepfig{\tSw}{a,b,c}{DegLow1,SPP-AC,DegLow3}{\tAc}{Degree of the low band: speedup per trace (left) and mean accuracy and coverage (right).}
\begin{itemize}
  \item Degree 1 ($-0.037\%$) and degree 3 ($+0.007\%$) are ties. Degree 1 has higher accuracy (54.3\%) and lower
        coverage (7.43\%), degree 3 the opposite.
  \item \textbf{Decision: keep 2.} Less than half of the prefetches are used here, so we prefetch less than in the
        middle band. The data does not show that 2 is better than 1 or 3.
\end{itemize}

\textbf{Bottom band (accuracy below 25\%)}\\
The degree is \textbf{1, never 0}. Without any prefetch no fill happens, the epoch never ends and the degree can never
recover. All four degrees are used on every trace (Table~\ref{tab:t1-epochs}), so every band is active.
\begin{table}[H]\centering\small
\begin{tabular}{@{}lccccc@{}}
\toprule
Trace & Degree 1 & Degree 2 & Degree 4 & Degree 6 & Total epochs \\
\midrule
Trace1 & 1  & 108 & 249 & 90  & 448 \\
Trace2 & 9  & 53  & 73  & 165 & 300 \\
Trace3 & 33 & 77  & 77  & 83  & 270 \\
\bottomrule
\end{tabular}
\caption{\code{SPP-Final}: number of 256-fill epochs that ended with each prefetch degree.}
\label{tab:t1-epochs}
\end{table}

\subsubsection{Values Fixed by Design (not swept)}
\begin{itemize}
  \item \textbf{Stride range $-64$ to $+63$:} a page has 64 lines, so larger strides cannot give a same-page prefetch.
  \item \textbf{History needed $=3$ strides:} equal to the signature length, so the PHT is used only with a full signature.
  \item \textbf{Minimum degree 1:} explained above.
  \item \textbf{Signature $=3$ strides of 4 bits, PHT of 4096 entries:} inherited from our first version. Three
        strides are enough for the assignment's example pattern. We did not test other lengths. A 4-bit stride mixes
        strides that differ by a multiple of 16.
  \item \textbf{PHT confidence (2 bits, 2 to start, 3 for chained strides):} chosen by reasoning (a chained prediction
        depends on the one before it), not by a sweep.
  \item \textbf{Start degree 3:} used only until the first epoch ends. Not tested.
  \item \textbf{Same-page rule, MSHR rule for the target level, uniform trigger:} taken from \code{IP-stride}.
  \item \textbf{Features added in groups:} \code{SPP-P1} changed three things at once (128 entries, 12-bit tag, valid bit)
        and \code{SPP-P2} two (stride range, history count). They were not tested one at a time. Their reasons are
        in Section~1.2, and the entries and tag were then swept above.
\end{itemize}

\subsubsection{Summary of Decisions}
\begin{table}[H]\centering\small
\begin{tabularx}{\textwidth}{@{}lllX@{}}
\toprule
Value & Chosen & Decided by & Reason \\
\midrule
Tracker entries              & 128                       & sweep + storage    & 256 is faster but about 22\% more storage for a small gain \\
IP tag                       & 8 bits                    & sweep              & faster and smaller than 12 bits \\
Prefetch filter              & 32 entries                & sweep + storage    & tie in speed, smaller \\
Epoch length                 & 256 fills                 & sweep + statistics & tie in speed; low sampling error \\
Thresholds                   & 75/50/25                  & sweep              & strict is slower, loose is less accurate \\
High-band degree             & 6                         & sweep              & 4 is worse, 8 is not clearly better \\
Middle-band degree           & 4                         & sweep (tie)        & kept as chosen \\
Low-band degree              & 2                         & sweep (tie)        & kept as chosen \\
Bottom-band degree           & 1                         & design             & degree 0 stops the epoch counter \\
Stride range, history needed & $-64..63$, 3              & design             & page size, signature length \\
Signature, PHT size          & 3 strides $\times$ 4 bits, 4096 & inherited    & not tested \\
PHT confidence limits        & 2-bit, 2 and 3            & design             & chained predictions need more confidence; not tested \\
Start degree                 & 3                         & design             & used only until the first epoch ends; not tested \\
\bottomrule
\end{tabularx}
\caption{Every value, how it was decided and why.}
\end{table}

\subsection{Final Choice and Storage}\label{sec:t1-final}
\code{SPP-Final} combines the two changes the sweep supports, the 8-bit tag and the 32-entry filter. Each was tested
alone, so we ran the combination as well (Tables~\ref{tab:t1-stages} and~\ref{tab:t1-avg}). Its geometric-mean
speedup is \textbf{1.01200} against 1.01079 for \code{SPP-AC} ($+0.120\%$), with mean accuracy 54.0\% against 52.6\%
and mean coverage 8.15\% against 7.92\%. Its per-trace speedups are 1.0098, 1.0133 and 1.0128 (\code{SPP-AC}:
1.0094, 1.0102, 1.0128).

\subsubsection{Storage Overhead of \code{SPP-Final}}
We assume 48-bit virtual addresses, so a cache-line address has $48-6 = 42$ bits. Only state needed in hardware is
counted (statistics counters are not).
\begin{table}[H]\centering\small
\begin{tabular}{@{}llr@{}}
\toprule
Field & Why this width & Bits \\
\midrule
\multicolumn{3}{@{}l}{\textit{IP tracker table, 128 entries}}\\
Valid             & entry in use                              & 1 \\
IP tag            & \code{ip \& 0xFF}                         & 8 \\
Last line address & previous cache line of this IP            & 42 \\
Last stride       & range $-64 \ldots +63$                    & 7 \\
Uniform flag      & only compared with 0, so 1 bit is enough  & 1 \\
Signature         & three 4-bit strides                       & 12 \\
History count     & 0 to 3                                    & 2 \\
LRU rank          & rank among 128 entries                    & 7 \\
\textbf{Entry total} &                                        & \textbf{80} \\
\textbf{Table total} & $128 \times 80$                        & \textbf{10240} \\
\midrule
\multicolumn{3}{@{}l}{\textit{Pattern History Table, 4096 entries, indexed by signature (no tag)}}\\
Predicted stride  & range $-64 \ldots +63$                    & 7 \\
Confidence        & 2-bit saturating counter (0 = empty)      & 2 \\
\textbf{Entry total} &                                        & \textbf{9} \\
\textbf{Table total} & $4096 \times 9$                        & \textbf{36864} \\
\midrule
\multicolumn{3}{@{}l}{\textit{Prefetch filter, 32 entries, direct-mapped}}\\
Valid             & line prefetched and not yet used          & 1 \\
Partial tag       & line address bits above the 5 index bits  & 10 \\
\textbf{Entry total} &                                        & \textbf{11} \\
\textbf{Table total} & $32 \times 11$                         & \textbf{352} \\
\midrule
\multicolumn{3}{@{}l}{\textit{Control registers}}\\
Fill counter $F$  & counts to $W = 256$                       & 8 \\
Useful counter $U$& at most $256 + 32 = 288$                  & 9 \\
Current degree    & values 1 to 6                             & 3 \\
\textbf{Total}    &                                           & \textbf{20} \\
\midrule
\textbf{Total}    & $10240 + 36864 + 352 + 20$                & \textbf{47476 bits} \\
                  &                                           & $= 5934.5$ bytes $\approx$ \textbf{5.8\,KiB} \\
\bottomrule
\end{tabular}
\caption{Bit-level storage of \code{SPP-Final}. \code{SPP-AC} (12-bit tag, 64-entry filter) needs 48340 bits, about 5.9\,KiB.}
\end{table}

\subsection{Final Results}
IPC, speedup, L1D load MPKI, prefetch accuracy and coverage of the \code{IP-stride} prefetcher and our submitted
prefetcher \code{SPP-Final}, on the three traces.
\begin{table}[H]\centering\small
\begin{tabular}{@{}lccccc@{}}
\toprule
Prefetcher & IPC & Speedup & L1D load MPKI & Accuracy (\%) & Coverage (\%) \\
\midrule
\grp{6}{Trace1}
\code{IP-stride} & 0.2334 & 1.0051 & 10.12 & 52.8 & 5.24 \\
\code{SPP-Final} & 0.2345 & 1.0098 & 9.68  & 57.8 & 9.37 \\
\midrule
\grp{6}{Trace2}
\code{IP-stride} & 0.2512 & 1.0002 & 13.44 & 38.1 & 0.73 \\
\code{SPP-Final} & 0.2545 & 1.0133 & 12.94 & 63.3 & 4.40 \\
\midrule
\grp{6}{Trace3}
\code{IP-stride} & 0.3830 & 1.0024 & 3.32  & 16.1 & 2.51 \\
\code{SPP-Final} & 0.3870 & 1.0128 & 3.04  & 41.0 & 10.69 \\
\bottomrule
\end{tabular}
\caption{IPC, speedup (over no prefetching), L1D load MPKI, prefetch accuracy and coverage of \code{IP-stride} and
the submitted prefetcher \code{SPP-Final}.}
\end{table}

\newcommand{\finopts}{cycle list name=final, bar width=9pt, enlarge x limits=0.25}
\begin{figure}[H]
\centering
\begin{subfigure}{0.49\textwidth}\centering
  \gbarplot{\tStageIPC}{ips,fin}{\finopts, ymin=0, ymax=0.45, ylabel={IPC},
    legend entries={IP-stride,SPP-Final}, legend columns=-1, legend to name=leg:final}
  \caption{IPC}
\end{subfigure}\hfill
\begin{subfigure}{0.49\textwidth}\centering
  \gbarplot{\tStageSpd}{ips,fin}{\finopts, ymin=1, ymax=1.016, ylabel={Speedup}}
  \caption{Speedup}
\end{subfigure}\\[4pt]
\begin{subfigure}{0.49\textwidth}\centering
  \gbarplot{\tStageMpki}{ips,fin}{\finopts, ymin=0, ymax=15, ylabel={L1D load MPKI}}
  \caption{L1D load MPKI}
\end{subfigure}\hfill
\begin{subfigure}{0.49\textwidth}\centering
  \gbarplot{\tStageAcc}{ips,fin}{\finopts, ymin=0, ymax=70, ylabel={Accuracy (\%)}}
  \caption{Prefetch accuracy}
\end{subfigure}\\[4pt]
\begin{subfigure}{0.49\textwidth}\centering
  \gbarplot{\tStageCov}{ips,fin}{\finopts, ymin=0, ymax=12, ylabel={Coverage (\%)}}
  \caption{Prefetch coverage}
\end{subfigure}\\[3pt]
\ref*{leg:final}
\caption{IPC, speedup, L1D load MPKI, prefetch accuracy and coverage of \code{IP-stride} and \code{SPP-Final} on the three traces.}
\end{figure}
